\section{Dataset y versiones}
\label{sec:dataset}

La fuente es el conjunto \href{https://www.kaggle.com/datasets/ajohrn/bikeshare-usage-in-london-and-taipei-network}{Bikeshare usage in London and Taipei} de Kaggle. \textbf{El alcance de este trabajo es únicamente la red de Londres}; los datos de Taipéi quedan fuera por decisión de equipo.

Se trabajó con el histórico de viajes en cinco versiones de tamaño creciente, más \texttt{london\_stations.csv}, un catálogo de 802 estaciones con nombre, latitud y longitud.

\begin{table}[htbp]
\centering
\caption{Versiones disponibles del dataset}
\label{tab:versiones}
\begin{tabular}{lrrrr}
\toprule
Versión & Archivo & Tamaño (GB) & Filas & Factor vs. XS \\
\midrule
XS & \texttt{london\_xs.csv} & 0.049 & 398,075 & $1.0\times$ \\
S  & \texttt{london\_s.csv}  & 0.488 & 3,972,436 & $10.0\times$ \\
M  & \texttt{london\_m.csv}  & 1.181 & 9,590,698 & $24.1\times$ \\
L  & \texttt{london\_l.csv}  & 2.363 & 19,148,881 & $48.1\times$ \\
XL & \texttt{london.csv}     & 4.725 & 38,215,560 & $96.0\times$ \\
\bottomrule
\end{tabular}
\end{table}

El esquema tiene nueve columnas: \texttt{rental\_id}, \texttt{duration}, \texttt{bike\_id}, \texttt{start\_rental\_date\_time}, \texttt{end\_rental\_date\_time}, \texttt{start\_station\_id}, \texttt{start\_station\_name}, \texttt{end\_station\_id}, \texttt{end\_station\_name}.

La exploración inicial, hecha con DuckDB en streaming para no cargar los archivos grandes en memoria, detectó cuatro problemas de calidad:

\begin{itemize}
  \item \texttt{end\_station\_id} tiene nulos, que corresponden a viajes sin devolución registrada. Se excluyen de las llegadas y de los pares origen-destino.
  \item \texttt{duration} contiene ceros y valores extremos de varios días. No intervienen en las tres preguntas, pero se documentan.
  \item \texttt{start\_station\_id} es entero mientras \texttt{end\_station\_id} se infiere como \texttt{DOUBLE} por causa de los nulos. Hay que castear explícitamente para que los identificadores comparen igual en las tres herramientas.
  \item El catálogo tiene 802 estaciones, pero Q2 y Q3 (Sección \ref{sec:respuestas}) reportan 808 identificadores de estación distintos, porque esa cifra sale directo de los \texttt{station\_id} que aparecen como origen o destino en los viajes, no del catálogo. Siete \texttt{station\_id} (391, 346, 304, 434, 241, 825, 823) no están en \texttt{london\_stations.csv}; la 391 por sí sola acumula 9,050 viajes. El join contra el catálogo es left join y las rellena con la etiqueta \texttt{(estación N)}, así que no rompen Q2 ni Q3, pero al no tener latitud ni longitud quedan fuera de la Figura \ref{fig:q2-mapa} (la estación 391 incluida), y su nombre se pierde de forma irrecuperable en la proyección a 4 columnas de la Fase 3 (Sección \ref{sec:fase3-diagnostico}).
\end{itemize}

El rango de fechas cambia con la versión: las versiones chicas cubren pocas semanas y las grandes cubren más meses. Eso altera la interpretación de la estacionalidad, no el método.
